\documentclass[10pt,a4paper]{article}
\usepackage[margin=2.2cm]{geometry}
\usepackage{fontspec}
\setmonofont{DejaVu Sans Mono}[Scale=0.84]
\newfontfamily\unifont{DejaVu Serif}[Scale=0.9]
\usepackage{amsmath,amssymb,mathtools}
\setcounter{MaxMatrixCols}{20}
\usepackage{longtable,booktabs,array}
\usepackage{fancyvrb}
\usepackage[dvipsnames]{xcolor}
\usepackage[colorlinks=true,linkcolor=NavyBlue,urlcolor=NavyBlue]{hyperref}
\usepackage{microtype}
\usepackage{needspace}
\usepackage{tocloft}\setlength{\cftbeforesecskip}{0.45em}
\setlength{\parindent}{0pt}\setlength{\parskip}{0.55em}
\setlength{\emergencystretch}{8em}
\makeatletter
\renewcommand\section{\@startsection{section}{1}{\z@}{-3.5ex \@plus -1ex \@minus -.2ex}{2.3ex \@plus.2ex}{\normalfont\Large\bfseries\raggedright}}
\renewcommand\subsection{\@startsection{subsection}{2}{\z@}{-3.25ex\@plus -1ex \@minus -.2ex}{1.5ex \@plus .2ex}{\normalfont\large\bfseries\raggedright}}
\renewcommand\subsubsection{\@startsection{subsubsection}{3}{\z@}{-3.25ex\@plus -1ex \@minus -.2ex}{1.5ex \@plus .2ex}{\normalfont\normalsize\bfseries\raggedright}}
\makeatother
\renewcommand{\arraystretch}{1.15}
\allowdisplaybreaks
\begin{document}
\title{The self-coupled dirac16complex in gravitational fields: symbols and their definitions\\[0.4em]\large A dictionary of every symbol used for the nine-channel self-coupling, the Fierz matrix, the covariant nonlinear field equation, chirality in a gravitational field and the Kohn–Sham theory in the three epochs — sections 15–16 of the notebook 7+1.nb}
\author{Repository: github.com/once-ere/7plus1 — prepared for Nash\\ Notebook: 7+1.nb (Wolfram Language 15.0.1)\\ Document generated by tools/make\_docs.py on 2026-10-28}
\date{}
\maketitle
\begin{abstract}
For each symbol this dictionary gives the mathematical notation, the name used in the Mathematica notebook 7+1.nb (where there is one), the defining equation and its meaning, grouped by topic in the order of introduction. Longer defining equations follow each table as numbered displayed formulas. Every equation needed to define a symbol is written in this document.
\end{abstract}
\tableofcontents
\Needspace{12\baselineskip}\section*{1. Conventions}\addcontentsline{toc}{section}{1. Conventions}
\begin{itemize}
\item Units: $\hbar=c=1$; a prime is $d/dx_0$.
\item Coordinates $X=(x_0,x_1,\dots,x_7)$ (\texttt{X = \{x0, …, x7\}}): $x_0$ time, $x_1,x_2,x_3$ ordinary space, $x_4,\dots,x_7$ auxiliary space.
\item Frame indices $A,B\in\lbrace0,\dots,7\rbrace$, raised and lowered with $\eta=\mathrm{diag}(+1,-1,\dots,-1)$ (\texttt{η17}), simply the squares of the octonion basis ($\eta_{AB}=\mathrm{Re}(e_Ae_B)$); coordinate indices $\mu,\nu\in\lbrace0,\dots,7\rbrace$; $i\in\lbrace1,2,3\rbrace$ ordinary, $k\in\lbrace4,\dots,7\rbrace$ auxiliary space. Repeated indices are summed.
\item Multi-indices $K,J$ are strictly increasing lists of frame indices; $|K|$ is the rank (the length of $K$); $K\in r$ means $|K|=r$.
\item Matrices: $1_n$ identity, $M^{\mathsf T}$ transpose, $M^\dagger$ conjugate transpose, $\mathrm{tr}$ trace, $[X,Y]=XY-YX$.
\item Mathematica names with Greek letters are written with named characters: \texttt{\textbackslash{}\allowbreak{}[CapitalTheta]up} is \texttt{Θup}, \texttt{\textbackslash{}\allowbreak{}[Eta]K} is \texttt{ηK}, \texttt{\textbackslash{}\allowbreak{}[Chi]17} is \texttt{χ17}, \texttt{\textbackslash{}\allowbreak{}[Omega]17} is \texttt{ω17}.
\item Signature: one time-like dimension ($x_0$) and seven space-like dimensions; the xAct list $\{7,1,0\}$ declares seven space-like dimensions and one time-like dimension, the matrix $\eta$ lists the squares of the octonion basis (below).
\end{itemize}
\textbf{The signature: $\{7,1,0\}$ and $\mathrm{diag}(+1,-1,\dots,-1)$.} The 7+1 spacetime has one time-like dimension ($x_0$, the direction of the octonion unit $e_0=1$) and seven space-like dimensions ($x_1,x_2,x_3$ ordinary space, $x_4,\dots,x_7$ auxiliary space). Two different objects describe this, and neither may be replaced by the other. (i) The tensor package xAct, used in section 12 of 7+1.nb exactly as in the author’s notebook Einstein-Lovelock-7+1-Nash.nb, declares the metric with \texttt{DefMetric[\{7, 1, 0\}, g71[-a, -b], CD, …]}: the list $\{7,1,0\}$ declares seven space-like dimensions, one time-like dimension and no degenerate dimension — the numbers of entries $+1$ and $-1$ of $\eta_{71}=\mathrm{diag}(-1,+1,\dots,+1)$ (\texttt{η71}), the same spacetime written in the sign convention of 4+4.nb (space-like $=+1$, time-like $=-1$). (ii) The matrix $\eta=-\eta_{71}=\mathrm{diag}(+1,-1,\dots,-1)$ (\texttt{η17}) is simply the squares of the octonion basis: $e_0^2=+1$ for the unit $e_0=1$ (time) and $e_i^2=-1$ for the seven imaginary units $e_1,\dots,e_7$ (space); more generally it is the real part of the product of two basis units, and it gives the quadratic form of the octonions:
\begin{gather*}
\eta_{AB}=\mathrm{Re}(e_Ae_B),\qquad e_0^2=+1,\qquad e_i^2=-1\ \ (i=1,\dots,7),\\ \mathrm{Re}(x\,x)=x_0^2-x_1^2-\cdots-x_7^2=\eta_{AB}\,x^Ax^B\qquad (x=x^Ae_A\in\mathbb O).
\end{gather*}
$\eta$ gives the signs of the Clifford relation $\Gamma_A\Gamma_B+\Gamma_B\Gamma_A=2\,\eta_{AB}\,1_{16}$ of the real 16×16 Dirac matrices: time squares to $+1_{16}$ and each of the seven space directions to $-1_{16}$, exactly like the octonion units (in the notation $Cl(p,q)$, with $p$ generators squaring to $+1$ and $q$ to $-1$, the algebra $Cl(1,7)\cong\mathbb R(16)$). It also gives the signs of the metric components: $g_{00}>0$ along the time-like direction, negative squares along the space-like directions. Every component computed by 7+1.nb follows from the component matrix of the metric; the declared list fixes for xAct the sign of the determinant of the metric, \texttt{SignDetOfMetric[g71]} $=-1$, which is the same in both sign conventions, since $\det \eta=\det \eta_{71}=-1$. 7+1.nb checks \texttt{SignatureOfMetric[g71]} $=\{7,1,0\}$, \texttt{SignDetOfMetric[g71]} $=-1$, $\mathrm{Re}(e_Ae_B)=\eta_{AB}$ and $\mathrm{Re}(x\,x)=\eta_{AB}x^Ax^B$.

\Needspace{12\baselineskip}\section*{2. Dirac matrices, chirality, and the field}\addcontentsline{toc}{section}{2. Dirac matrices, chirality, and the field}
{\small
\begin{longtable}{>{\raggedright\arraybackslash}p{\dimexpr 0.1500\linewidth-2\tabcolsep\relax}>{\raggedright\arraybackslash}p{\dimexpr 0.2100\linewidth-2\tabcolsep\relax}>{\raggedright\arraybackslash}p{\dimexpr 0.3600\linewidth-2\tabcolsep\relax}>{\raggedright\arraybackslash}p{\dimexpr 0.2800\linewidth-2\tabcolsep\relax}}
\toprule
\textbf{symbol} & \textbf{Mathematica} & \textbf{definition} & \textbf{meaning} \\
\midrule\endfirsthead
\toprule
\textbf{symbol} & \textbf{Mathematica} & \textbf{definition} & \textbf{meaning} \\
\midrule\endhead
\endfoot
\bottomrule\endlastfoot
$\Gamma_A$ & \texttt{Γ17[A]} & (2.1) & real 16×16 Dirac matrices of 7+1 spacetime \\
$\Gamma^A$ & \texttt{Γ17up[A]} & $\eta^{AA}\Gamma_A$ (no sum) &  \\
$\tau^O_i$ & \texttt{τO[i]} & real antisymmetric 8×8, $(\tau^O_i)^2=-1_8$, anticommuting & the seven imaginary octonion units as matrices \\
$\omega$ & \texttt{ω17} & $\Gamma_0\Gamma_1\cdots\Gamma_7$ & volume element; $\omega^2=-1_{16}$ \\
$\chi$ & \texttt{χ17} & $i\omega$ & chirality; $\chi^2=1_{16}$, Hermitian \\
$P_{L,R}$ & \texttt{PL17}, \texttt{PR17} & $\tfrac12(1_{16}\mp\chi)$ & chiral projections \\
$E_L$, $E_R$ & \texttt{EL}, \texttt{ER} & (2.2) & chiral basis (16×8) \\
$\phi_L$, $\phi_R$ & \texttt{phiLc}, \texttt{phiRc} & $E_{L,R}^\dagger\psi$ & 8-component chiral parts \\
$\psi$ & \texttt{dirac16complex}; \texttt{psN}, \texttt{psR} & 16 complex anticommuting components & the field; in the checks random Gaussian-integer spinors \\
$\bar\psi$ & \texttt{pbS}, \texttt{pbR} & $\psi^\dagger\Gamma_0$ (independent symbols in the checks) & Dirac adjoint \\
$\alpha_j$, $\beta$ & \texttt{αD[j]}, \texttt{βD} & $\Gamma^0\Gamma^j$, $\Gamma_0$ & Hamiltonian Dirac matrices \\
$m$ & \texttt{mD}; \texttt{mks} &  & mass of dirac16complex ($m=1$ in section 16) \\*
\end{longtable}}
\begin{align*}
\Gamma_0&=\begin{pmatrix}0&1_8\\1_8&0\end{pmatrix},\qquad \Gamma_i=\begin{pmatrix}\tau^O_i&0\\0&-\tau^O_i\end{pmatrix},\qquad \Gamma_A\Gamma_B+\Gamma_B\Gamma_A=2\eta_{AB}1_{16},&&(2.1) \\
E_L&=\frac1{\sqrt2}\begin{pmatrix}1_8\\-i\,1_8\end{pmatrix},\qquad E_R=\frac1{\sqrt2}\begin{pmatrix}-i\,1_8\\1_8\end{pmatrix},\qquad \chi E_L=-E_L,\quad \chi E_R=E_R .&&(2.2)
\end{align*}
\Needspace{12\baselineskip}\section*{3. Gravity}\addcontentsline{toc}{section}{3. Gravity}
{\small
\begin{longtable}{>{\raggedright\arraybackslash}p{\dimexpr 0.1500\linewidth-2\tabcolsep\relax}>{\raggedright\arraybackslash}p{\dimexpr 0.2100\linewidth-2\tabcolsep\relax}>{\raggedright\arraybackslash}p{\dimexpr 0.3600\linewidth-2\tabcolsep\relax}>{\raggedright\arraybackslash}p{\dimexpr 0.2800\linewidth-2\tabcolsep\relax}}
\toprule
\textbf{symbol} & \textbf{Mathematica} & \textbf{definition} & \textbf{meaning} \\
\midrule\endfirsthead
\toprule
\textbf{symbol} & \textbf{Mathematica} & \textbf{definition} & \textbf{meaning} \\
\midrule\endhead
\endfoot
\bottomrule\endlastfoot
$e^A{}_\mu$, $e_A{}^\mu$ & \texttt{fd["e"]}, \texttt{fd["ei"]} & $g_{\mu\nu}=\eta_{AB}e^A{}_\mu e^B{}_\nu$, $e_A{}^\mu e^A{}_\nu=\delta^\mu_\nu$ & frame and inverse frame \\
$\Gamma^\mu(X)$ & \texttt{fdCos["GamUp"]} & $e_A{}^\mu\Gamma^A$ & generalized frame matrices \\
$\{{}^{\lambda}_{\mu\nu}\}$ & \texttt{cv["chr"]} & (3.1) & Christoffel symbols (never written $\Gamma$) \\
$\omega_\mu{}^A{}_B$ & \texttt{fd["om"]} & (3.2) & canonical spin connection \\
$\Omega_\mu$ & \texttt{fdCos["Om"]} & (3.2) & spinor connection \\
$D_\mu$ & — & $\partial_\mu+\Omega_\mu$ & covariant derivative of spinors \\
$a$, $b$ & \texttt{a3[x0]}, \texttt{b4[x0]}; \texttt{aE}, \texttt{bE} & (3.3) & scale factors of ordinary and auxiliary space \\
$H_a$, $H_b$ & \texttt{Ha}, \texttt{Hb}; \texttt{hE}, \texttt{kE} & $a'/a$, $b'/b$ & Hubble rates \\
$H_a'$, $H_b'$ & \texttt{dHa}, \texttt{dHb} & $dH_a/dx_0$, $dH_b/dx_0$ &  \\
$E^{(k)\mu}{}_\nu$ & \texttt{ECos[[k+1]]} & (3.4) & Lovelock tensors ($E^{(1)}$ = Einstein tensor) \\
$\alpha_k$ & \texttt{alf[k]} & $\alpha_1=1$ & Lovelock couplings ($\alpha_2$, $\alpha_3$) \\
$\bar\Lambda$ & \texttt{lamb} &  & cosmological constant \\
$\kappa$ & \texttt{kap} & $8\pi G_8$ & gravitational coupling \\
$p$, $q$ & \texttt{pKs}, \texttt{qKs} & $p=(1+2\sqrt2)/7$, $q=(2-3\sqrt2)/14$ & Einstein–Kasner exponents: $a=x_0^p$, $b=x_0^q$ \\*
\end{longtable}}
\begin{align*}
\{{}^{\lambda}_{\mu\nu}\}&=\tfrac12g^{\lambda\rho}\big(\partial_\mu g_{\rho\nu}+\partial_\nu g_{\rho\mu}-\partial_\rho g_{\mu\nu}\big),&&(3.1) \\
\omega_\mu{}^A{}_B&=e^A{}_\nu\big(\partial_\mu e_B{}^\nu+\{{}^{\nu}_{\mu\lambda}\}e_B{}^\lambda\big),\qquad \Omega_\mu=\tfrac14\eta_{AC}\,\omega_\mu{}^C{}_B\,\Gamma^A\Gamma^B,&&(3.2) \\
ds^2&=dx_0^2-a^2\big(dx_1^2+dx_2^2+dx_3^2\big)-b^2\big(dx_4^2+\cdots+dx_7^2\big),&&(3.3) \\
E^{(k)\mu}{}_\nu&=-\frac{1}{2^{k+1}}\,\delta^{\mu\alpha_1\beta_1\cdots\alpha_k\beta_k}_{\nu\gamma_1\delta_1\cdots\gamma_k\delta_k}R^{\gamma_1\delta_1}{}_{\alpha_1\beta_1}\cdots R^{\gamma_k\delta_k}{}_{\alpha_k\beta_k},\quad \delta^{\mu_1\cdots\mu_n}_{\nu_1\cdots\nu_n}=\det\big[\delta^{\mu_i}_{\nu_j}\big].&&(3.4)
\end{align*}
\Needspace{12\baselineskip}\section*{4. The nine channels and the Fierz matrix}\addcontentsline{toc}{section}{4. The nine channels and the Fierz matrix}
{\small
\begin{longtable}{>{\raggedright\arraybackslash}p{\dimexpr 0.1500\linewidth-2\tabcolsep\relax}>{\raggedright\arraybackslash}p{\dimexpr 0.2100\linewidth-2\tabcolsep\relax}>{\raggedright\arraybackslash}p{\dimexpr 0.3600\linewidth-2\tabcolsep\relax}>{\raggedright\arraybackslash}p{\dimexpr 0.2800\linewidth-2\tabcolsep\relax}}
\toprule
\textbf{symbol} & \textbf{Mathematica} & \textbf{definition} & \textbf{meaning} \\
\midrule\endfirsthead
\toprule
\textbf{symbol} & \textbf{Mathematica} & \textbf{definition} & \textbf{meaning} \\
\midrule\endhead
\endfoot
\bottomrule\endlastfoot
$K$ & \texttt{mIdx} & the 256 increasing subsets of $\lbrace0,\dots,7\rbrace$ & multi-indices \\
$r=|K|$ & \texttt{rk[K]} & length of $K$ & rank, $0\le r\le8$ \\
$\Gamma^K$ & — & $\Gamma^{A_1}\cdots\Gamma^{A_r}$, $\Gamma^{()}=1_{16}$ & product of Dirac matrices \\
$\Theta^K$ & \texttt{Θup[K]} & (4.1) & Hermitian-phased products; $\bar\psi\Theta^K\psi$ real \\
$\eta_K$ & \texttt{ηK[K]} & $\eta_{A_1A_1}\cdots\eta_{A_rA_r}$ & sign of the multi-index \\
$\Theta_K$ & \texttt{Θdn[K]} & $\eta_K\Theta^K$ & index lowered; $\Theta_K\Theta^K=1_{16}$ \\
$I_r$ & \texttt{Ichan[r, pb, ps]} & (4.2) & Lorentz invariant of rank $r$ \\
$g_r$ & \texttt{gg[r]}, \texttt{gSymb} &  & the nine couplings, $r=0,\dots,8$ \\
$g_S,g_V,g_A,g_P$ & \texttt{gS}, \texttt{gV}, \texttt{gA}, \texttt{gP} & $g_0=g_S$, $g_1=g_V$, $g_7=-g_A$, $g_8=-g_P$ & the four couplings of section 8 of 7+1.nb \\
$\mathcal L_{\rm self}$ & \texttt{LselfS} & $-\tfrac12\sum_rg_rI_r$ & self-coupling Lagrangian \\
$S$ & \texttt{sLL} & $\exp\big(\tfrac18w_{AB}[\Gamma^A,\Gamma^B]\big)$, $w_{AB}=-w_{BA}$ real & local Lorentz transformation of spinors \\
$F_{rs}$ & \texttt{fierzF} & (4.3) & 9×9 Fierz matrix, $F^2=1$ \\
$R$ & \texttt{RrF} & a 16×16 matrix & density matrix (random in the checks) \\
$D_s$ & \texttt{dS9[R]} & $\sum_{J\in s}\mathrm{tr}(\Theta_JR)\,\mathrm{tr}(\Theta^JR)$ & direct (Hartree) invariants \\
$G_s$ & \texttt{Geff} & $G=(1-F^{\mathsf T})g$ & Hartree–Fock (Fierz-transformed) couplings \\
$G_S$, $G_V$ & \texttt{gSks}, \texttt{gVks} & $G_S=G_0$, $G_V=G_1$ & couplings of uniform matter ($-100000$, $10000$) \\
$\rho_S$, $\rho$ & \texttt{rhoSs}, \texttt{rhos} & $\bar\psi\psi$, $\psi^\dagger\psi$ (expectation values) & scalar and number densities \\*
\end{longtable}}
\begin{align*}
\Theta^K&=i^{\sigma_r}\,\Gamma^{A_1}\cdots\Gamma^{A_r},\qquad \sigma_r=\tfrac{r(r-1)}{2}\bmod 2,&&(4.1) \\
I_r&=\sum_{|K|=r}(\bar\psi\Theta_K\psi)(\bar\psi\Theta^K\psi),&&(4.2) \\
F_{rs}&=\frac1{16}\sum_l(-1)^{rs-l}\binom{s}{l}\binom{8-s}{r-l}=\frac1{256}\sum_{K\in r}\mathrm{tr}\big(\Theta_K\Theta_J\Theta^K\Theta^J\big)\ \ (|J|=s).&&(4.3)
\end{align*}
The Fierz identity for commuting spinors $a,b,c,d$:
\begin{equation*}
\sum_{K\in r}(a\Theta_Kb)(c\Theta^Kd)=\sum_sF_{rs}\sum_{J\in s}(a\Theta_Jd)(c\Theta^Jb).
\end{equation*}
\Needspace{12\baselineskip}\section*{5. Field equation and Kohn–Sham quantities}\addcontentsline{toc}{section}{5. Field equation and Kohn–Sham quantities}
{\small
\begin{longtable}{>{\raggedright\arraybackslash}p{\dimexpr 0.1500\linewidth-2\tabcolsep\relax}>{\raggedright\arraybackslash}p{\dimexpr 0.2100\linewidth-2\tabcolsep\relax}>{\raggedright\arraybackslash}p{\dimexpr 0.3600\linewidth-2\tabcolsep\relax}>{\raggedright\arraybackslash}p{\dimexpr 0.2800\linewidth-2\tabcolsep\relax}}
\toprule
\textbf{symbol} & \textbf{Mathematica} & \textbf{definition} & \textbf{meaning} \\
\midrule\endfirsthead
\toprule
\textbf{symbol} & \textbf{Mathematica} & \textbf{definition} & \textbf{meaning} \\
\midrule\endhead
\endfoot
\bottomrule\endlastfoot
$\Sigma(X)$ & — & (5.1) & Kohn–Sham self-energy (direct minus exchange) \\
$\varphi_n$ & — & $(i\Gamma^\mu D_\mu-m-\Sigma)\varphi_n=0$ & Kohn–Sham orbitals \\
$\mathcal R$ & — & $\sum_{n\ \rm occ}\varphi_n\bar\varphi_n$ & Kohn–Sham density matrix \\
$k=(\vec p,\vec q)$ & — & comoving momentum & $\vec p$ ordinary, $\vec q$ auxiliary \\
$k_F$ & \texttt{kF} & radius of the comoving Fermi ball & Fermi momentum \\
$k_{F,\rm sat}$ & \texttt{kFsat} & $p=0$ at $a=b=1$ & saturation Fermi momentum \\
$\theta$ & — & $|k_a|=k\cos\theta$, $|k_b|=k\sin\theta$ & angle between ordinary and auxiliary momentum \\
$\omega$ (energy) & \texttt{om} & $\sqrt{k^2\cos^2\theta/a^2+k^2\sin^2\theta/b^2+M^2}$ & single-particle energy (not the volume element) \\
$M$, $m^*$ & \texttt{Mx}, \texttt{ksGap} & $M=m+G_S\rho_S(M)$ & effective mass (gap equation) \\
$V_0$ & — & $G_V\rho$ & vector potential \\
$\langle f\rangle$ & \texttt{ksMoments} & (5.2) & average over the comoving Fermi ball \\
$\varepsilon$, $p_a$, $p_b$ & \texttt{ksStress} & (5.3) & energy density, pressures along ordinary / auxiliary space \\
$\Pi$ & \texttt{pint} & (5.3) & interaction pressure \\
$\varepsilon_F$ & — & $G_V\rho+\sqrt{k_F^2+M^2}$ & Fermi energy (flat spacetime) \\
$E/A$ & — & $\varepsilon/\rho$ & energy per particle \\
$T^\mu{}_\nu$ & — & $\mathrm{diag}(\varepsilon,-p_a,-p_a,-p_a,-p_b,\dots,-p_b)$ & energy–momentum tensor of the Kohn–Sham matter \\
— & \texttt{evolveEL[α2, α3, κ, Ha0]} & returns $\lbrace H_b(0),\lbrace a,b,H_a,H_b\rbrace(20),$ violation$\rbrace$ & self-consistent Einstein–Lovelock + Kohn–Sham evolution \\*
\end{longtable}}
\begin{align*}
\Sigma&=\sum_rg_r\sum_{|K|=r}\big[\mathrm{tr}(\Theta_K\mathcal R)\,\Theta^K-\Theta^K\mathcal R\,\Theta_K\big],&&(5.1) \\
\langle f\rangle&=\frac{1}{2\pi^4a^3b^4}\int_0^{k_F}dk\,k^6\int_0^{\pi/2}d\theta\,\cos^2\theta\sin^3\theta\,f,\qquad \rho=\langle1\rangle,\quad \rho_S=\langle M/\omega\rangle,&&(5.2) \\
\varepsilon&=\langle\omega\rangle-\frac{(M-m)^2}{2G_S}+\frac{G_V}2\rho^2,\quad \Pi=\frac{(M-m)^2}{2G_S}+\frac{G_V}2\rho^2,&&(5.3) \\
p_a&=\Big\langle\frac{k^2\cos^2\theta}{3a^2\omega}\Big\rangle+\Pi,\quad p_b=\Big\langle\frac{k^2\sin^2\theta}{4b^2\omega}\Big\rangle+\Pi .
\end{align*}
\Needspace{12\baselineskip}\section*{6. Symbols that look alike}\addcontentsline{toc}{section}{6. Symbols that look alike}
{\small
\begin{longtable}{>{\raggedright\arraybackslash}p{\dimexpr 0.4200\linewidth-2\tabcolsep\relax}>{\raggedright\arraybackslash}p{\dimexpr 0.5800\linewidth-2\tabcolsep\relax}}
\toprule
\textbf{symbol} & \textbf{not to be confused with} \\
\midrule\endfirsthead
\toprule
\textbf{symbol} & \textbf{not to be confused with} \\
\midrule\endhead
\endfoot
\bottomrule\endlastfoot
$\{7,1,0\}$ (xAct signature list: seven space-like dimensions, one time-like) & $\eta=\mathrm{diag}(+1,-1,\dots,-1)$ (simply the squares of the octonion basis, $\eta_{AB}=\mathrm{Re}(e_Ae_B)$) \\
$\omega=\Gamma_0\cdots\Gamma_7$ (volume element) & $\omega$ (single-particle energy), $\omega_\mu{}^A{}_B$ (spin connection) \\
$\chi$ (chirality) & $\chi_k$ (rescaled mode, $\psi=(a^3b^4)^{-1/2}\chi_k$) \\
$\Theta^K$ (upper multi-index) & $\Theta_K=\eta_K\Theta^K$ (lower multi-index) \\
$g_r$ (bare couplings) & $G_s$ (Fierz-transformed couplings), $g_{\mu\nu}$ (metric) \\
$F_{rs}$ (Fierz matrix) & $\mathcal F_{\mu\nu}$ (spin curvature, not used here) \\
$R$ (density matrix) & $\mathcal R$ (Kohn–Sham density matrix), $R^{\gamma\delta}{}_{\alpha\beta}$ (Riemann tensor) \\
$p$ (Kasner exponent; pressure in flat spacetime) & $p_a$, $p_b$ (pressures), $\vec p$ (ordinary momentum) \\
$q$ (Kasner exponent) & $\vec q$ (auxiliary momentum) \\
$\alpha_k$ (Lovelock couplings) & $\alpha_j$ (Dirac matrices $\Gamma^0\Gamma^j$) \\
$\kappa=8\pi G_8$ & — \\*
\end{longtable}}
\end{document}
